\documentclass[addpoints,spanish, 12pt]{exam}

\usepackage[default,osfigures,scale=0.95]{opensans} %% Alternatively
%% use the option '' instead of 'default' to replace the
%% sans serif font only.
\usepackage[T1]{fontenc}

\usepackage[spanish]{babel}
\usepackage{graphicx}
\usepackage{epstopdf}
\usepackage{float}

\epstopdfsetup{outdir=./} 
\selectlanguage{spanish}
\usepackage[utf8]{inputenc}

\usepackage{tikz}
\usetikzlibrary{scopes}


\author{Felipe Antonio Riveroll Aguirre}
\pagestyle{headandfoot}
\runningheadrule
\firstpageheader{}
                {Instituto Cristóbal Colón A.C.\\
                Sexto Exámen Parcial de Física\\
                Marzo 2014}
                {}

\runningheader{Física}
{Sexto examen parcial}
{Marzo 13, 2014}
\firstpagefooter{}{{\tiny Elaborado por Felipe Riveroll Aguirre}}{}
\runningfooter{}{{\tiny Elaborado por Felipe Riveroll Aguirre}}{Página \thepage\ de \numpages}
\pointpoints{Punto}{Puntos}
\bonuspointpoints{punto extra}{puntos extra}
\bracketedpoints
\setlength\answerlinelength{0.5in}
\usepackage{letltxmacro,xpatch}
\xpatchcmd{\answerline}% <cmd>
  {\par\nobreak\vskip\answerskip}% <search>
  {\leaders\hbox{\hskip0.2em .\hskip0.2em}\hskip 0pt plus 1filll}% <replace>
  {}{}% <success><failure>
\xpatchcmd{\answerline}{\fi \par}{\fi}{}{}% Remove line break after \answerline
\makeatletter
\LetLtxMacro{\oldanswerline}{\answerline}
\RenewDocumentCommand{\answerline}{s o}{%
  \begingroup
  \IfBooleanTF{#1}
    {\def\@queslevel{\relax}}% \answerline*
    {}% \answerline
  \IfNoValueTF{#2}
    {\oldanswerline[{}]}% \answerline
    {\oldanswerline[#2]}% \answerline[..]
  \endgroup
}
\makeatother


\begin{document}
\vspace{0.1in}
\makebox[\textwidth]{Nombre:\enspace\hrulefill}


\begin{questions}
\question[25]
Complete con (V)erdadero o (F)also según corresponda

\begin{parts}
\part[5] La energía cinética se conserva tanto en el choque elástico como en el inelástico  \answerline*[F]
\part[5] El momento lineal de un objeto es proporcional a la masa del objeto \answerline*[V] 
\part[5] El cambio en el momento desarrollado por un cuerpo es igual al impulso que actúa en él \answerline*[V]
\part[5] En una colisión elástica los objetos chocan y se adhieren entre sí\answerline*[F]
\part[5] En una colisión inelástica se conserva la energía cinética \answerline*[F] 
\end{parts}
\question[25]
Liste en orden de mayor a menor momento $(p_{x})_{1}$ a $(p_{x})_{5}$

\begin{figure}[H]
\centering
\includegraphics[width=7in]{masas1}
\end{figure}

\question[25]
Un cuerpo que tiene una masa de 10 kg se deja caer al suelo y justo antes de chocar contra el suelo tiene 500 J de energía cinética ¿Desde que altura se dejó caer?

\question[25] Una esfera metálica de 50 kg es arrojada desde la azotea de un edificio, cuando se encuentra a 20 m del nivel del suelo, ¿Cuál será su velocidad final a esa altura?

\bonusquestion[50] Un carro de 1500 Kg se mueve hacia el norte conforme a la gráfica velocidad-tiempo mostrada en la siguiente figura

\begin{parts}
\part[10] Determine el cambio de la energía cinética durante los primeros 7 segundos.
\part[10] Para Determinar que tan lejos ha viajado este carro en 7 segundos, no pueden ser utilizadas las 3 ecuaciones básicas de cinemática ($v_{f}=v_{i}+at$, $\Delta x=v_{i}t + \frac{1}{2} at^{2}$, $v^2_f = v^2_i + 2a \Delta x$). Explique la razón para este caso.
\part[10] Utilize la gráfica de velocidad-tiempo para estimar la distancia que el carro ha viajado en 7 segundos.
\part[10] ¿Cuál fué el trabajo neto realizado durante su recorrido de 7 segundos? 
\part[10] Determine la potencia necesaria para que el carro desarrolle este movimiento. Exprese la potencia en caballos de fuerza ($1 hp = 0.746 kW$)
\end{parts}
\end{questions}


\end{document}